\chapter{Operator Schmidt rank at cuts inside a bulk region}
\label{ch:alvv_interior_cut_rank}

Item~4 of Theorem~4.2 of Arad, Landau, Vazirani and
Vidick~\cite{Arad2016RigorousRG} bounds the operator Schmidt rank of a
polynomial filter at cuts that lie inside one of the three bulk regions,
away from the two windows. At such a cut every term of the Hamiltonian
except one bulk operator lies on one side, so the Hamiltonian is a rank-one
operator on each side plus the bulk operator.

\begin{lemma}[Rank at a cut inside one bulk region]
    \label{thm:alvv_bulk_cut_polynomial_rank}
    \lean{TNLean.MPS.AreaLaw.operatorSchmidtRank_aeval_sum_le_pow_of_bulk}
    \leanok
    Let $S$ be a set of sites of a finite system with local dimension $q$,
    and let $H=\sum_\beta F_\beta$ be a finite sum of operators. Suppose that
    every $F_\beta$ with $\beta\ne\beta_0$ acts on $S$ or on its complement,
    and that $F_{\beta_0}$ has operator Schmidt rank at most $B$ across
    $S$ and its complement. Then for every complex polynomial $P$ of degree
    at most $k$,
    \begin{align}
        \operatorname{OSR}_{S\mid S^c}\bigl(P(H)\bigr)
         & \le(B+2)^{k+1}.
        \label{eq:alvv_bulk_cut_polynomial_rank}
    \end{align}
\end{lemma}
\begin{proof}\leanok
    Collect the terms acting on $S$ into $A$ and those acting on $S^c$ into
    $C$. In the coordinates of $S$ and $S^c$, $A=A'\otimes1$ and
    $C=1\otimes C'$, so each has operator Schmidt rank at most one, and
    subadditivity gives rank at most $B+2$ for $H=A+F_{\beta_0}+C$.
    Submultiplicativity gives rank at most $(B+2)^j$ for $H^j$, including
    $H^0=1$. Since $P(H)$ is a combination of $H^j$ with $j\le k$,
    \begin{align}
        \operatorname{OSR}_{S\mid S^c}\bigl(P(H)\bigr)
         & \le\sum_{j=0}^k(B+2)^j\le(B+2)^{k+1},\notag
    \end{align}
    where the last step uses $B+2\ge2$.
\end{proof}

\begin{lemma}[Rank at cuts inside the three bulk regions]
    \label{thm:alvv_two_window_interior_cut_rank}
    \lean{TNLean.MPS.AreaLaw.operatorSchmidtRank_aeval_two_window_middle_cut_le,
        TNLean.MPS.AreaLaw.operatorSchmidtRank_aeval_two_window_left_cut_le,
        TNLean.MPS.AreaLaw.operatorSchmidtRank_aeval_two_window_right_cut_le}
    \leanok
    \uses{thm:alvv_bulk_cut_polynomial_rank}
    Let $L,q,\ell,a,b$ be natural numbers with $a\le b$, and let $F_\alpha$,
    $\alpha\in\mathcal I$, be the two-window family of
    Lemma~\ref{thm:chain_two_window_polynomial_rank}: the window terms
    $F_{x_r}$ and $F_{y_r}$ supported on $[a+r,a+r+2)$ and $[b+r,b+r+2)$
    for $0\le r<2\ell$, and the bulk operators $F_\lambda,F_\mu,F_\rho$
    supported on $J_L=[0,a+1)$, $J_M=[a+2\ell,b+1)$ and
    $J_R=[b+2\ell,L)$. Put $H=\sum_\alpha F_\alpha$, let $P$ have degree at
    most $k$, and consider the cut between the sites $c-1$ and $c$, that is,
    the bipartition $[0,c)\mid[c,L)$. Suppose that one of the following holds:
    \begin{enumerate}
        \item $a+2\ell+1\le c\le b$, and $F_\mu$ has operator Schmidt rank at
              most $B$ at the cut;
        \item $c\le a$, and $F_\lambda$ has operator Schmidt rank at most $B$
              at the cut;
        \item $b+2\ell+1\le c$, and $F_\rho$ has operator Schmidt rank at
              most $B$ at the cut.
    \end{enumerate}
    Then
    \begin{align}
        \operatorname{OSR}_{[0,c)\mid[c,L)}\bigl(P(H)\bigr)
         & \le(B+2)^{k+1}.
        \label{eq:alvv_two_window_interior_cut_rank}
    \end{align}
    This is item~4 of Theorem~4.2 of~\cite{Arad2016RigorousRG}, where $B$
    bounds the Schmidt rank of the bulk operator of the region containing
    the cut. For $B\ge2$ the right side of
    (\ref{eq:alvv_two_window_interior_cut_rank}) is at most $B^{2k+2}$,
    the printed form $B^{O(k)}$. For $B\le1$ the printed form fails: if
    $F_\mu=0$ and the parts of $H$ on the two sides of the cut are diagonal
    with distinct entries $\alpha_i$ and $\beta_j$, the rank of $P(H)$ at the
    cut is the rank of the matrix $\bigl(P(\alpha_i+\beta_j)\bigr)_{ij}$,
    which is $k+1$ for suitable entries. The two boundary terms cannot be
    absorbed into the base, and
    (\ref{eq:alvv_two_window_interior_cut_rank}) is the corrected statement.
\end{lemma}
\begin{proof}\leanok
    \uses{thm:alvv_bulk_cut_polynomial_rank}
    In each case every term other than the named bulk operator lies on one
    side of the cut. In the first case the left window, which ends at site
    $a+2\ell$, and $J_L$ lie in $[0,c)$, while the right window and $J_R$ lie
    in $[c,L)$. In the second case all window terms and $J_M,J_R$ lie in
    $[c,L)$. In the third case all window terms and $J_L,J_M$ lie in $[0,c)$.
    Lemma~\ref{thm:alvv_bulk_cut_polynomial_rank} gives
    (\ref{eq:alvv_two_window_interior_cut_rank}).
\end{proof}
